\chapter{NKCode : l'éditeur du dépôt}

Vous pouvez écrire du Nkentseu dans n'importe quel éditeur de texte. Mais le
dépôt en fournit un, écrit avec le moteur lui-même, et qui connaît Jenga : c'est
\texttt{NKCode}. Ce chapitre le prend en main --- assez pour que vous n'ayez plus
à quitter l'éditeur pour compiler, chercher ou déboguer.

\begin{nkretenir}
NKCode est \textbf{écrit avec NKCanvas et NKGui}, les deux modules que ce cours
enseigne. C'est donc aussi le plus gros exemple à votre disposition : 56 fichiers
d'une application réelle. Quand une question du cours vous restera, ouvrez son
code --- il répond souvent.
\end{nkretenir}

\section{Le lancer}

\begin{nkcode}{Commandes -- a taper dans un terminal a la racine du depot}
jenga build --target NKCode --config Release
.\Build\Bin\Release-Windows\NKCode\NKCode.exe
\end{nkcode}

\begin{nkpiege}
\textbf{Un NKCode fraîchement construit ne démarre pas sous Windows}, et le build
ne le dit pas. Il meurt sur \texttt{python312.dll est introuvable} : quatre DLL
du runtime Python embarqué sont posées par \texttt{scripts/MakeNkCodeDist.py},
pas par \texttt{jenga build}.

Pour un essai rapide, préfixez le \texttt{PATH} sans rien copier :

\begin{nkcode}{Contournement}
$env:PATH = "<depot>\Externals\Libs\PythonEmbed\runtime;$env:PATH"
.\Build\Bin\Release-Windows\NKCode\NKCode.exe
\end{nkcode}

Le diagnostic qui l'a trouvé : comparer le dossier de sortie avec celui d'un
poste où NKCode tourne --- il contient quatre fichiers de plus. \emph{Un build
vert et un exécutable mort ne sont pas contradictoires} : ils mesurent deux
choses différentes.
\end{nkpiege}

\section{Les raccourcis}

\begin{center}
\begin{tabular}{@{}ll@{}}
\toprule
\textbf{Raccourci} & \textbf{Action} \\
\midrule
\texttt{Ctrl+P} & palette de commandes --- \emph{commencez par là} \\
\texttt{Ctrl+N} & nouveau fichier \\
\texttt{Ctrl+S} & enregistrer \\
\texttt{Ctrl+F} & rechercher (non modal) \\
\texttt{Ctrl+G} & aller à la ligne \\
\texttt{Ctrl+D} & dupliquer la ligne ou la sélection \\
\texttt{Ctrl+L} & sélectionner la ligne (s'étend si répété) \\
\texttt{F2} & renommer l'identifiant \\
\texttt{F9} & poser un point d'arrêt \\
\texttt{F5} & déboguer --- lance \texttt{jenga gdb}, points d'arrêt transmis \\
\texttt{F1} & documentation \\
\texttt{Ctrl+Q} & quitter \\
\bottomrule
\end{tabular}
\end{center}

\begin{nkretenir}
\textbf{\texttt{Ctrl+P} rend tous les autres facultatifs.} La palette liste les
commandes par leur nom : si vous ne vous rappelez pas d'un raccourci, tapez le
verbe. C'est aussi la façon de \emph{découvrir} ce que l'éditeur sait faire.
\end{nkretenir}

\section{IntelliSense : clangd}

NKCode parle le protocole LSP et s'appuie sur \texttt{clangd}. Vous obtenez donc
les diagnostics en temps réel, le survol qui montre un type, et l'aller-à-la
définition --- les mêmes que dans un IDE commercial, parce que c'est le même
moteur d'analyse.

\begin{nkpiege}
\texttt{clangd} a besoin de savoir \emph{comment} chaque fichier est compilé :
quels chemins d'inclusion, quelles macros. Sans cette information, il souligne en
rouge des inclusions parfaitement valides.

\textbf{Un fichier souligné en rouge qui compile parfaitement n'est pas un défaut
du code : c'est un défaut de configuration de l'analyseur.} Ne « corrigez » rien
avant d'avoir lancé un vrai build --- c'est \texttt{jenga} qui a raison, pas le
soulignement.
\end{nkpiege}

\section{Construire sans quitter l'éditeur}

NKCode connaît Jenga. Le point d'arrêt que vous posez avec \texttt{F9} est
transmis à \texttt{jenga gdb} par \texttt{F5} : vous n'écrivez aucune ligne de
commande de débogage.

\begin{nkretenir}
\textbf{Ce que le débogueur vous donne, et que l'affichage ne donne pas} :
l'\emph{état} au moment de l'arrêt --- toutes les variables, la pile d'appels,
qui a appelé qui.

Un \texttt{printf} répond à une question que vous avez déjà formulée. Un point
d'arrêt répond à celles que vous n'aviez pas pensé à poser. Quand vous ne
comprenez pas ce qui se passe, le débogueur est plus rapide, toujours.
\end{nkretenir}

\section{Les autres services}

\begin{center}
\begin{tabular}{@{}ll@{}}
\toprule
\textbf{Service} & \textbf{Ce qu'il fait} \\
\midrule
Recherche / remplacement & \texttt{Ctrl+F}, \texttt{Ctrl+H}, non modaux \\
Multilingue & 8 langues, changement sans redémarrer \\
Assistant IA & discussion, contexte du fichier, commandes \\
Émulateurs & lance Android et HarmonyOS depuis l'éditeur \\
\bottomrule
\end{tabular}
\end{center}

\begin{nkretenir}
« Non modal » veut dire que la barre de recherche ne bloque pas l'édition : vous
tapez dans le fichier pendant qu'elle est ouverte. C'est un détail qui change
l'usage --- une boîte de dialogue modale interrompt le geste, une barre non
modale l'accompagne.
\end{nkretenir}

\begin{nkpratique}
Ouvrez le dépôt dans NKCode et faites les cinq gestes qui reviennent tous les
jours :

\begin{enumerate}
\item \texttt{Ctrl+P}, puis tapez « ouvrir » --- lisez les commandes proposées ;
\item ouvrez \texttt{Kernel/Runtime/NKEvent/src/NKEvent/NkPointerEvent.h} et
      survolez \texttt{NkPointerPhase} : le type doit s'afficher ;
\item posez un point d'arrêt (\texttt{F9}) dans le \texttt{nkmain} d'une petite
      application, lancez \texttt{F5}, et regardez la pile d'appels ;
\item cherchez \texttt{NkCanvasApp} dans tout le dépôt --- combien de fichiers ?
\item changez la langue de l'interface, et vérifiez que rien ne redémarre.
\end{enumerate}
\end{nkpratique}

\begin{nkbilan}
NKCode est un confort, jamais une obligation : tout ce que fait ce livre se fait
aussi avec un éditeur quelconque et un terminal. Ce qu'il apporte est
concentré --- la palette de commandes, l'auto-complétion réelle par clangd, la
construction sans quitter l'éditeur. Retenez surtout que son IntelliSense lit le
\emph{même} projet que le compilateur : s'il ne trouve pas un en-tête, c'est
souvent le projet qui est mal décrit, pas l'éditeur qui se trompe.
\end{nkbilan}

\begin{nkquestions}
\begin{enumerate}
\item Que lit clangd pour savoir où sont vos en-têtes ?
\item Que fait \texttt{Ctrl+P}, et en quoi diffère-t-il d'une recherche de
      fichier ordinaire ?
\item Pourquoi une barre de recherche \emph{non modale} change-t-elle l'usage ?
\item Peut-on construire un projet sans quitter l'éditeur ? Par quel moyen ?
\item NKCode est-il nécessaire pour suivre ce cours ? Justifiez.
\end{enumerate}
\end{nkquestions}

\begin{nkdemo}
Prouvez que l'auto-complétion lit bien votre projet, et non une supposition.

Créez un en-tête à vous, dans un dossier que le projet ne déclare pas.
Incluez-le depuis un fichier source : l'auto-complétion ne doit rien proposer.
Ajoutez ensuite le dossier aux chemins d'inclusion du \texttt{.jenga},
rechargez, et vérifiez qu'elle propose désormais vos symboles.

\textbf{Ce que la démonstration établit} : que le service dépend d'une
\emph{donnée} --- la description du projet --- et non d'une magie de l'éditeur.
C'est ce qui vous dira quoi corriger le jour où il ne trouve plus rien.
\end{nkdemo}

\begin{nklogique}
Un collègue vous dit : « l'auto-complétion ne marche pas, NKCode est cassé ».

Énumérez au moins quatre causes possibles, en distinguant celles qui viennent de
l'éditeur de celles qui viennent du projet. Pour chacune, donnez la
\emph{première chose à regarder} --- et non le correctif.

Puis dites laquelle est la plus probable, et pourquoi c'est justement celle
qu'on soupçonne en dernier.
\end{nklogique}
